\usepackage{luatexja} % LuaTeXで日本語を使うためのパッケージ
\usepackage{luatexja-fontspec} % LuaTeX用の日本語フォント設定
% lualatex main && upmendex -r -c -s main.ist -g main && lualatex main  

% --- 数学関連 ---
\usepackage{amsmath, amssymb, amsfonts, mathtools, bm, amsthm} % 基本的な数学パッケージ
\usepackage{type1cm, upgreek} % 数式フォントとギリシャ文字k
\usepackage{physics, mhchem} % 物理や化学の記号や式の表記を簡単にする

% --- 表関連 ---
\usepackage{multirow, longtable, tabularx, array, colortbl, dcolumn, diagbox} % 表のレイアウトを柔軟にする
\usepackage{tablefootnote, truthtable} % 表中に注釈を追加、真理値表
\usepackage{tabularray} % 高度な表組みレイアウト

% --- グラフィック関連 ---
\usepackage{tikz, graphicx} % 図の描画と画像の挿入
% \usepackage{background} % ウォーターマークの設定
\usepackage{caption, subcaption} % 図や表のキャプション設定
\usepackage{float, here} % 図や表の位置指定

% --- レイアウトとページ設定 ---
\usepackage{fancyhdr} % ページヘッダー、フッター、余白の設定
\usepackage[top = 20truemm, bottom = 20truemm, left = 20truemm, right = 20truemm]{geometry}
\usepackage{fancybox, ascmac} % ボックスのデザイン

% --- 色とスタイル ---
\usepackage{xcolor, color, colortbl, tcolorbox} % 色とカラーボックス
\usepackage{listings, jvlisting} % コードの色付けとフォーマット

% --- 参考文献関連 ---
% \usepackage[style=phys, backend=biber]{biblatex}
\usepackage[backend=biber, sortcites=true, style=phys]{biblatex}
\usepackage{usebib} % 参考文献の管理と挿入
\usepackage{url, hyperref} % URLとリンクの設定

% --- その他の便利なパッケージ ---
\usepackage{footmisc} % 脚注のカスタマイズ
\usepackage{multicol} % 複数段組
\usepackage{comment} % コメントアウトの拡張
\usepackage{siunitx} % 単位の表記
\usepackage{docmute, needspace}
\usepackage{makeidx}

% \renewbibmacro*{title}{%
%   \printtext[title]{\printfield[titlecase]{title}.\addspace}%
% }
% --- 定理スタイルと数式設定 ---
\theoremstyle{definition} % 定義スタイル
\numberwithin{equation}{section} % 式番号をサブセクション単位でリセット

% --- hyperrefの設定 ---
\hypersetup{
  setpagesize = false,
  bookmarks = true,
  bookmarksdepth = tocdepth,
  bookmarksnumbered = true,
  colorlinks = false,
  pdftitle = {}, % PDFタイトル
  pdfsubject = {}, % PDFサブジェクト
  pdfauthor = {}, % PDF作者
  pdfkeywords = {} % PDFキーワード
}

% --- siunitxの設定 ---
\sisetup{
  table-format = 1.5, % 小数点以下の桁数
  table-number-alignment = center, % 数値の中央揃え
}


% --- その他の設定 ---
\allowdisplaybreaks % 数式の途中改ページ許可
\newcolumntype{t}{!{\vrule width 0.1pt}} % 新しいカラムタイプ
\newcolumntype{b}{!{\vrule width 1.5pt}} % 太いカラム
\UseTblrLibrary{amsmath, booktabs, counter, diagbox, functional, hook, html, nameref, siunitx, varwidth, zref} % tabularrayのライブラリ
\setlength{\columnseprule}{0.4pt} % カラム区切り線の太さ
% \captionsetup[figure]{font = bf} % 図のキャプションの太字設定
% \captionsetup[table]{font = bf} % 表のキャプションの太字設定
% \captionsetup[lstlisting]{font = bf} % コードのキャプションの太字設定
% \captionsetup[subfigure]{font = bf, labelformat = simple} % サブ図のキャプション設定
\setcounter{secnumdepth}{4} % セクションの深さ設定
\newcolumntype{d}{D{.}{.}{5}} % 数値のカラム
\newcolumntype{M}[1]{>{\centering\arraybackslash}m{#1}} % センター揃えのカラム
\DeclareMathOperator{\diag}{diag}
\everymath{\displaystyle} % 数式のスタイル

\setcounter{tocdepth}{0} % 目次の深さ
\makeatletter
\@addtoreset{equation}{chapter} % セクションごとに式番号をリセット
\makeatother

\newcommand{\braketg}{\braket{\r{g}}}
\ExecuteBibliographyOptions{
  sorting=none, biblabel=brackets, articletitle=true, chaptertitle=false,
  pageranges=false, maxnames=99, minnames=1, hyperref=auto, abbreviate=true,
  date=year, arxiv=abs, isbn=false, url=false, doi=true, eprint=true, giveninits=true
}
\DeclareSourcemap{
    \maps[datatype=bibtex, overwrite=true]{
        \map{
            \step[
                fieldsource=journal,
                match=\regexp{Physical\sReview\sLetters},
                replace={Phys. Rev. Lett.}
            ]
            \step[
                fieldsource=journal,
                match=\regexp{Physical\sReview\sA},
                replace={Phys. Rev. A}
            ]
            \step[
                fieldsource=journal,
                match=\regexp{Physical\sReview\sB},
                replace={Phys. Rev. B}
            ]
            \step[
                fieldsource=journal,
                match=\regexp{Physical\sReview\sC},
                replace={Phys. Rev. C}
            ]
            \step[
                fieldsource=journal,
                match=\regexp{Physical\sReview\sD},
                replace={Phys. Rev. D}
            ]
            \step[
                fieldsource=journal,
                match=\regexp{Physical\sReview\sE},
                replace={Phys. Rev. E}
            ]
            \step[
                fieldsource=journal,
                match=\regexp{Physical\sReview\sX},
                replace={Phys. Rev. X}
            ]
            \step[
                fieldsource=journal,
                match=\regexp{Reviews\sof\sModern\sPhysics},
                replace={Rev. Mod. Phys.}
            ]
            \step[
                fieldsource=journal,
                match=\regexp{Journal\sof\sApplied\sPhysics},
                replace={J. Appl. Phys.}
            ]
            \step[
                fieldsource=journal,
                match=\regexp{Applied\sPhysics\sA},
                replace={Appl. Phys. A}
            ]
            \step[
                fieldsource=journal,
                match=\regexp{Applied\sPhysics\sB},
                replace={Appl. Phys. B}
            ]
            \step[
                fieldsource=journal,
                match=\regexp{Applied\sPhysics\sLetters},
                replace={Appl. Phys. Lett.}
            ]
            \step[
                fieldsource=journal,
                match=\regexp{Review\sof\sScientific\sInstruments},
                replace={Rev. Sci. Instrum.}
            ]
            \step[
                fieldsource=journal,
                match=\regexp{Nature},
                replace={Nature}
            ]
            \step[
                fieldsource=journal,
                match=\regexp{Science},
                replace={Science}
            ]
            \step[
                fieldsource=journal,
                match=\regexp{Nature\sPhysics},
                replace={Nat. Phys.}
            ]
            \step[
                fieldsource=journal,
                match=\regexp{Nature\sCommunications},
                replace={Nat. Commun.}
            ]
            \step[
                fieldsource=journal,
                match=\regexp{Scientific\sReports},
                replace={Sci. Rep.}
            ]
            \step[
                fieldsource=journal,
                match=\regexp{Journal\sof\sPhysics\sA:\sMathematical\sand\sTheoretical},
                replace={J. Phys. A: Math. Theor.}
            ]
            \step[
                fieldsource=journal,
                match=\regexp{Journal\sof\sPhysics\sB:\sAtomic,\sMolecular\sand\sOptical\sPhysics},
                replace={J. Phys. B: At. Mol. Opt. Phys.}
            ]
            \step[
                fieldsource=journal,
                match=\regexp{Journal\sof\sPhysics:\sCondensed\sMatter},
                replace={J. Phys.: Condens. Matter}
            ]
            \step[
                fieldsource=journal,
                match=\regexp{Annalen\sder\sPhysik},
                replace={Ann. Phys.}
            ]
            \step[
                fieldsource=journal,
                match=\regexp{Europhysics\sLetters},
                replace={Europhys. Lett.}
            ]
            \step[
                fieldsource=journal,
                match=\regexp{New\sJournal\sof\sPhysics},
                replace={New J. Phys.}
            ]
            \step[
                fieldsource=journal,
                match=\regexp{Physics\sLetters\sA},
                replace={Phys. Lett. A}
            ]
            \step[
                fieldsource=journal,
                match=\regexp{Physics\sLetters\sB},
                replace={Phys. Lett. B}
            ]
            \step[
                fieldsource=journal,
                match=\regexp{Nuclear\sPhysics\sA},
                replace={Nucl. Phys. A}
            ]
            \step[
                fieldsource=journal,
                match=\regexp{Nuclear\sPhysics\sB},
                replace={Nucl. Phys. B}
            ]
            \step[
                fieldsource=journal,
                match=\regexp{Journal\sof\sHigh\sEnergy\sPhysics},
                replace={J. High Energy Phys.}
            ]
            \step[
                fieldsource=journal,
                match=\regexp{Classical\sand\sQuantum\sGravity},
                replace={Class. Quantum Grav.}
            ]
            \step[
                fieldsource=journal,
                match=\regexp{Monthly\sNotices\sof\sthe\sRoyal\sAstronomical\sSociety},
                replace={Mon. Not. R. Astron. Soc.}
            ]
            \step[
                fieldsource=journal,
                match=\regexp{IEEE\sTransactions\son\sAntennas\sand\sPropagation},
                replace={IEEE Trans. Antennas Propag.}
            ]
            \step[
                fieldsource=journal,
                match=\regexp{IEEE\sTransactions\son\sCommunications},
                replace={IEEE Trans. Commun.}
            ]
            \step[
                fieldsource=journal,
                match=\regexp{IEEE\sTransactions\son\sSignal\sProcessing},
                replace={IEEE Trans. Signal Process.}
            ]
            \step[
                fieldsource=journal,
                match=\regexp{IEEE\sTransactions\son\sPhotonics},
                replace={IEEE Trans. Photonics}
            ]
            \step[
                fieldsource=journal,
                match=\regexp{IEEE\sJournal\son\sQuantum\sElectronics},
                replace={IEEE J. Quantum Electron.}
            ]
            \step[
                fieldsource=journal,
                match=\regexp{IEEE\sJournal\son\sSelected\sTopics\sin\sQuantum\sElectronics},
                replace={IEEE J. Sel. Top. Quantum Electron.}
            ]
            \step[
                fieldsource=journal,
                match=\regexp{IEEE\sTransactions\son\sWireless\sCommunications},
                replace={IEEE Trans. Wireless Commun.}
            ]
            \step[
                fieldsource=journal,
                match=\regexp{IEEE\sTransactions\son\sMicrowave\sTheory\sand\sTechniques},
                replace={IEEE Trans. Microwave Theory Tech.}
            ]
            \step[
                fieldsource=journal,
                match=\regexp{IEEE\sTransactions\son\sInstrumentation\sand\sMeasurement},
                replace={IEEE Trans. Instrum. Meas.}
            ]
            \step[
                fieldsource=journal,
                match=\regexp{IEEE\sTransactions\son\sElectromagnetic\sCompatibility},
                replace={IEEE Trans. Electromagn. Compat.}
            ]
            \step[
                fieldsource=journal,
                match=\regexp{IEEE\sTransactions\son\sImage\sProcessing},
                replace={IEEE Trans. Image Process.}
            ]
            \step[
                fieldsource=journal,
                match=\regexp{IEEE\sTransactions\son\sCircuit\sand\sSystem\sII:\sExpress\sBriefs},
                replace={IEEE Trans. Circuits Syst. II: Express Briefs}
            ]
            \step[
                fieldsource=journal,
                match=\regexp{IEEE\sAccess},
                replace={IEEE Access}
            ]
            \step[
                fieldsource=journal,
                match=\regexp{IEEE\sTransactions\son\sNanotechnology},
                replace={IEEE Trans. Nanotechnol.}
            ]
            \step[
                fieldsource=journal,
                match=\regexp{IEEE\sTransactions\son\sFuzzy\sSystems},
                replace={IEEE Trans. Fuzzy Syst.}
            ]
            \step[
                fieldsource=journal,
                match=\regexp{IEEE\sTransactions\son\sControl\sSystems\stheory\sand\sApplications},
                replace={IEEE Trans. Control Syst. Theory Appl.}
            ]
            \step[
                fieldsource=journal,
                match=\regexp{IEEE\sTransactions\son\sSmart\sGrid},
                replace={IEEE Trans. Smart Grid}
            ]
            \step[
                fieldsource=journal,
                match=\regexp{IEEE\sTransactions\son\sVehicular\sTechnology},
                replace={IEEE Trans. Veh. Technol.}
            ]
            \step[
                fieldsource=journal,
                match=\regexp{IEEE\sTransactions\son\sPower\sElectronics},
                replace={IEEE Trans. Power Electron.}
            ]
            \step[
                fieldsource=journal,
                match=\regexp{IEEE\sTransactions\son\sIndustrial\sElectronics},
                replace={IEEE Trans. Ind. Electron.}
            ]
            \step[
                fieldsource=journal,
                match=\regexp{IEEE\sTransactions\son\sComputers},
                replace={IEEE Trans. Comput.}
            ]
            \step[
                fieldsource=journal,
                match=\regexp{IEEE\sTransactions\son\sSoftware\sEngineering},
                replace={IEEE Trans. Softw. Eng.}
            ]
        }
    }
}
\newcommand{\inner}[2]{\left\langle #1, #2 \right\rangle}
\renewcommand{\figurename}{図}
\renewcommand{\i}{\mathrm{i}} % 複素数単位i
\renewcommand{\laplacian}{\grad^2} % ラプラシアンの記号
\renewcommand{\thesubfigure}{(\alph{subfigure})} % サブ図の番号形式
\newcommand{\m}[3]{\multicolumn{#1}{#2}{#3}} % マルチカラムのショートカット
\renewcommand{\r}[1]{\mathrm{#1}} % mathrmのショートカット
\newcommand{\e}{\mathrm{e}} % 自然対数の底e
\newcommand{\Ef}{E_{\mathrm{F}}} % フェルミエネルギー
\renewcommand{\c}{\si{\degreeCelsius}} % 摂氏記号
\renewcommand{\d}{\r{d}} % d記号
\renewcommand{\t}[1]{\texttt{#1}} % タイプライタフォント
\newcommand{\kb}{k_{\mathrm{B}}} % ボルツマン定数
\renewcommand{\epsilon}{\varepsilon}
\newcommand{\fullref}[1]{\textbf{\ref{#1} \nameref{#1}}}
\newcommand{\reff}[1]{{図\ref{#1}}} % 図参照のショートカット
\newcommand{\reft}[1]{{表\ref{#1}}} % 表参照のショートカット
\newcommand{\refe}[1]{{式\eqref{#1}}} % 式参照のショートカット
\newcommand{\refp}[1]{{コード\ref{#1}}} % コード参照のショートカット
\renewcommand{\lstlistingname}{コード} % コードリストの名前
\renewcommand{\theequation}{\thesection.\arabic{equation}} % 式番号の形式
\renewcommand{\footrulewidth}{0.4pt} % フッターの線
\newcommand{\mar}[1]{\textcircled{\scriptsize #1}} % 丸囲み文字
\newcommand{\combination}[2]{{}_{#1} \mathrm{C}_{#2}} % 組み合わせ
\newcommand{\thline}{\noalign{\hrule height 0.1pt}} % 細い横線
\newcommand{\bhline}{\noalign{\hrule height 1.5pt}} % 太い横線
\renewcommand{\appendixname}{付録}


\bibliography{ref.bib}

% \DeclareFieldFormat{eprint}{arXiv:\href{http://arxiv.org/abs/#1}{#1}}
\DeclareFieldFormat{eprint}{arXiv:\href{http://arxiv.org/abs/ #1}{#1} (\thefield{year})}
% 出力のカスタマイズ（タイトルを省略し、arXiv形式を適用）
% \DeclareFieldFormat[article]{title}{\mkbibemph{#1}}  % 必要に応じて
\AtEveryBibitem{\clearfield{doi}}  % DOIフィールドを消す
\AtEveryBibitem{\clearfield{url}}  % URLフィールドを消す
\DeclareFieldFormat[article]{note}{\mkbibparens{arXiv:#1}}  % arXivの表示形式を調整
% \AtEveryBibitem{
%   \clearfield{title}
% }
% --- メタ情報 ---
\title{論文まとめノート}
\date{更新日: \today}
\author{Haruki AOKI}
